% !PRINTED 312
\weberpage{312}
%<<CONT
\weq{3}{(1+y)^n f\left(\frac{\alpha+\beta y}{1+y}\right) = \varphi(y) = b_0\,y^n + b_1\,y^{n-1} + \cdots + b_n,}
worin die $b_0$, $b_1 \ldots b_n$ ganze homogene Functionen $n^{\mathrm{ten}}$ Grades von $\alpha$ und $\beta$ und lineare Functionen der Co\"efficienten von $f(x)$ sind. Wenn nun $x$ von $\alpha$ bis $\beta$ wächst, so geht $y$, gleichfalls wachsend, von $0$ bis $\infty$. So viele Wurzeln also $f(x)$ zwischen $\alpha$ und $\beta$ hat, genau so viele positive Wurzeln hat $\varphi(y)$.

Es folgt also aus dem \gesperrt{Cartesischen} Satz:

\gesperrt{IX. Die Anzahl der Wurzeln von $f(x)$ zwischen $\alpha$ und $\beta$ ist so gross oder um eine gerade Zahl kleiner, wie die Anzahl der Zeichenwechsel in der Reihe}
\[
b_0,\ b_1,\ b_2 \ldots b_n.
\]

Während also bei dem \gesperrt{Budan'schen} Theorem die Vorzeichenwechsel in zwei Werthreihen abzuzählen und zu vergleichen sind, verlangt das \gesperrt{Jacobi'sche} Kriterium nur die Abzählung der Vorzeichenwechsel in einer Werthreihe.

Auf dieselbe Weise lässt sich auch die \gesperrt{Newton'sche} Regel umgestalten, wenn man die Grössen
\[
B_\nu = \nu(n-\nu)\,b_\nu^2 - (\nu+1)(n-\nu+1)\,b_{\nu-1}\,b_{\nu+1}
\]
in Betracht zieht, und man erhält so den Satz:

\gesperrt{X. Die Anzahl der Wurzeln von $f(x)$ zwischen $\alpha$ und $\beta$ ist so gross oder um eine gerade Zahl kleiner, wie die Zahl der Wechselfolgen in der Doppelreihe}
\[
\begin{aligned}
&b_0,\ b_1,\ b_2 \ldots b_n,\\
&B_0,\ B_1,\ B_2 \ldots B_n.
\end{aligned}
\]

\begin{center}
\S.\,103.

Klein's geometrische Vergleichung der verschiedenen Kriterien.
\end{center}

Um die Tragweite der verschiedenen Kriterien unter einander zu vergleichen, hat \gesperrt{Klein} eine geometrische Interpretation angewandt, die wir jetzt besprechen wollen. Wenn man nicht mit Räumen von mehr Dimensionen operiren und damit die geometrische Anschauung verlieren will, muss man sich hierbei
%CONT>>
